% Options for packages loaded elsewhere
\PassOptionsToPackage{unicode}{hyperref}
\PassOptionsToPackage{hyphens}{url}
\documentclass[
  10pt,
  a4paper,
]{article}
\usepackage{xcolor}
\usepackage[margin=22mm]{geometry}
\usepackage{amsmath,amssymb}
\setcounter{secnumdepth}{-\maxdimen} % remove section numbering
\usepackage{iftex}
\ifPDFTeX
  \usepackage[T1]{fontenc}
  \usepackage[utf8]{inputenc}
  \usepackage{textcomp} % provide euro and other symbols
\else % if luatex or xetex
  \usepackage{unicode-math} % this also loads fontspec
  \defaultfontfeatures{Scale=MatchLowercase}
  \defaultfontfeatures[\rmfamily]{Ligatures=TeX,Scale=1}
\fi
\usepackage{lmodern}
\ifPDFTeX\else
  % xetex/luatex font selection
\fi
% Use upquote if available, for straight quotes in verbatim environments
\IfFileExists{upquote.sty}{\usepackage{upquote}}{}
\IfFileExists{microtype.sty}{% use microtype if available
  \usepackage[]{microtype}
  \UseMicrotypeSet[protrusion]{basicmath} % disable protrusion for tt fonts
}{}
\makeatletter
\@ifundefined{KOMAClassName}{% if non-KOMA class
  \IfFileExists{parskip.sty}{%
    \usepackage{parskip}
  }{% else
    \setlength{\parindent}{0pt}
    \setlength{\parskip}{6pt plus 2pt minus 1pt}}
}{% if KOMA class
  \KOMAoptions{parskip=half}}
\makeatother
\setlength{\emergencystretch}{3em} % prevent overfull lines
\providecommand{\tightlist}{%
  \setlength{\itemsep}{0pt}\setlength{\parskip}{0pt}}
\usepackage{mathrsfs}
\usepackage{mathtools}
\usepackage{xurl}
\emergencystretch=3em
\setcounter{tocdepth}{1}
\newcommand{\Beta}{\mathrm{B}}
\DeclareUnicodeCharacter{00B2}{\ensuremath{^2}}
\DeclareUnicodeCharacter{00B1}{\ensuremath{\pm}}
\DeclareUnicodeCharacter{00D7}{\ensuremath{\times}}
\DeclareUnicodeCharacter{2192}{\ensuremath{\to}}
\DeclareUnicodeCharacter{21D2}{\ensuremath{\Rightarrow}}
\DeclareUnicodeCharacter{21A6}{\ensuremath{\mapsto}}
\DeclareUnicodeCharacter{2208}{\ensuremath{\in}}
\DeclareUnicodeCharacter{2209}{\ensuremath{\notin}}
\DeclareUnicodeCharacter{2212}{\ensuremath{-}}
\DeclareUnicodeCharacter{2227}{\ensuremath{\wedge}}
\DeclareUnicodeCharacter{2228}{\ensuremath{\vee}}
\DeclareUnicodeCharacter{2260}{\ensuremath{\ne}}
\DeclareUnicodeCharacter{2264}{\ensuremath{\le}}
\DeclareUnicodeCharacter{2265}{\ensuremath{\ge}}
\DeclareUnicodeCharacter{2297}{\ensuremath{\otimes}}
\DeclareUnicodeCharacter{00B3}{\ensuremath{^3}}
\DeclareUnicodeCharacter{2074}{\ensuremath{^4}}
\DeclareUnicodeCharacter{2076}{\ensuremath{^6}}
\DeclareUnicodeCharacter{03BB}{\ensuremath{\lambda}}
\DeclareUnicodeCharacter{03BC}{\ensuremath{\mu}}
\DeclareUnicodeCharacter{03B5}{\ensuremath{\epsilon}}
\DeclareUnicodeCharacter{03C4}{\ensuremath{\tau}}
\DeclareUnicodeCharacter{03A6}{\ensuremath{\Phi}}
\DeclareUnicodeCharacter{03B1}{\ensuremath{\alpha}}
\DeclareUnicodeCharacter{03C3}{\ensuremath{\sigma}}
\DeclareUnicodeCharacter{03B4}{\ensuremath{\delta}}
\DeclareUnicodeCharacter{03C0}{\ensuremath{\pi}}
\DeclareUnicodeCharacter{039B}{\ensuremath{\Lambda}}
\DeclareUnicodeCharacter{2202}{\ensuremath{\partial}}
\DeclareUnicodeCharacter{0393}{\ensuremath{\Gamma}}
\DeclareUnicodeCharacter{2205}{\ensuremath{\varnothing}}
\DeclareUnicodeCharacter{221E}{\ensuremath{\infty}}
\DeclareUnicodeCharacter{2200}{\ensuremath{\forall}}
\DeclareUnicodeCharacter{00B9}{\ensuremath{^1}}
\DeclareUnicodeCharacter{00B0}{\ensuremath{{}^\circ}}
\DeclareUnicodeCharacter{211D}{\ensuremath{\mathbb{R}}}
\usepackage{bookmark}
\IfFileExists{xurl.sty}{\usepackage{xurl}}{} % add URL line breaks if available
\urlstyle{same}
\hypersetup{
  hidelinks,
  pdfcreator={LaTeX via pandoc}}

\author{}
\date{}

\begin{document}

{
\setcounter{tocdepth}{3}
\tableofcontents
}
\section{One fixed primitive test and the full translation
criterion}\label{one-fixed-primitive-test-and-the-full-translation-criterion}

This proof constructs one specified compactly supported smooth test,
chosen independently of every zeta zero. Its Mellin transform is nonzero
at every possible off-critical zero. The full reflected Weil pairing of
that test with its translates is bounded on the positive half-line
exactly when RH holds. The proof does not establish that boundedness.

The original minus-exponent Mellin convention, actual zero multiset,
boundary filter, and global Laplace argument are those of
\href{BOUNDARY_PRIMITIVE_ALL_ZERO_RECEIVER.tex}{BPZ1--BPZ23}. That
entire TeX source was read for this derivation. The explicit
construction here makes its detecting test independent of a selected or
hypothetical zero. The method is an application of the classical
Weil-positivity mechanism, with no assertion of historical priority for
a bounded-translation criterion.

\subsection{1. A fixed infinite
convolution}\label{a-fixed-infinite-convolution}

The bump used here is classical. Juan Arias de Reyna, \emph{An
infinitely differentiable function with compact support: Definition and
properties}, Rev.~Real Acad. Ciencias Madrid \textbf{76} (1982), 21--38,
author's English translation
\href{https://arxiv.org/abs/1702.05442v1}{arXiv:1702.05442v1}, Theorem 1
and equations (1), (4), (9), constructs the function \(\varphi\) used in
the following exact comparison. His Fourier convention is
\(\widehat\varphi(\zeta)=\int\varphi(t)e^{-2\pi i t\zeta}\,dt\). In the
unchanged coordinates of this proof, \(b_r(v)=r^{-1}\varphi(v/r)\) and
\(G_r(z)=\widehat\varphi(rz/(2\pi i))\). Indeed, substituting \(v=rt\)
in the integral gives that transform, and its factor indexed by
\(h\ge0\) is \(\sinh(rz2^{-h-1})/(rz2^{-h-1})\), the factor indexed by
\(j=h+1\) in (UP3). Fourier uniqueness proves the density identity. The
factor \(r^{-1}\), the support \([-r,r]\), and total mass one are
retained. The complete direct construction below is kept for the reader;
it is not a new invention of this bump. Arias de Reyna credits Jessen
and Wintner (1935) for its Fourier-transform construction; that original
article was not inspected for this citation revision.

Throughout the construction retain \[
r=\frac1{64},\qquad a_j=r2^{-j}\quad(j\ge1),\qquad
\sum_{j\ge1}a_j=r,\quad \sum_{j\ge1}a_j^2=\frac{r^2}{3}.
\tag{UP1}
\] Let \(u_j(v)=(2a_j)^{-1}\mathbf1_{[-a_j,a_j]}(v)\), and let
\(b_{r,N}=u_1*\cdots*u_N\) be its finite convolution density. This is a
nonnegative probability density, is real and even, and is supported in
\([-r_N,r_N]\), where \(r_N=r(1-2^{-N})\). These assertions follow
successively from the integral formula for convolution: its integral is
the product of the integrals, convolution adds support intervals, and
reflection preserves the convolution of even functions.

Use the Fourier convention
\(\widehat b(y)=\int_{\mathbb R}b(v)e^{-iyv}\,dv\). Direct integration
on each interval gives \[
\widehat b_{r,N}(y)=\prod_{j=1}^N\frac{\sin(a_jy)}{a_jy},\qquad
G_{r,N}(z)=\int_{\mathbb R}b_{r,N}(v)e^{-zv}\,dv
=\prod_{j=1}^N\frac{\sinh(a_jz)}{a_jz},
\tag{UP2}
\] with each quotient assigned its removable value 1 at zero. In
particular \(\widehat b_{r,N}(y)=G_{r,N}(iy)\).

The entire functions in (UP2) converge uniformly on compact subsets of
\(\mathbb C\) to an entire function \[
G_r(z)=\prod_{j=1}^\infty\frac{\sinh(a_jz)}{a_jz},
\qquad G_r(0)=1.
\tag{UP3}
\] Here is the required product estimate. The power series for
\(\sinh w/w\) gives \[
\left|\frac{\sinh w}{w}-1\right|
\le\frac{|w|^2e^{|w|}}6.
\tag{UP4}
\] For the coefficient comparison, \((2k+1)!\ge6(2k-2)!\) for \(k\ge1\),
so the tail is bounded by \(|w|^2/6\) times a subseries of \(e^{|w|}\).
On \(|z|\le R\), (UP4) and (UP1) bound the sum of the deviations of the
factors from 1 by \(R^2e^{rR/2}r^2/18\). The tail of this convergent sum
is uniform on that disc. The estimate
\(\prod(1+|d_j|)\le\exp(\sum|d_j|)\) then shows that the product partial
sums are uniformly Cauchy. Their locally uniform limit is entire,
proving (UP3).

For real \(y\), each factor in \(\widehat b_{r,N}\) has modulus at most
1. For every fixed integer \(M\ge1\), every \(N\ge M\), and \(|y|\ge1\),
keeping its first \(M\) factors gives the explicit bound \[
|\widehat b_{r,N}(y)|
\le r^{-M}2^{M(M+1)/2}|y|^{-M}.
\tag{UP5}
\] For \(|y|\le1\) the bound 1 holds. The same two bounds hold for
\(G_r(iy)\) by taking the limit. Thus the limiting Fourier function is
rapidly decreasing, and every polynomial multiple of it is integrable.

Define \[
b_r(v)=\frac1{2\pi}\int_{\mathbb R}G_r(iy)e^{iyv}\,dy.
\tag{UP6}
\] For every integer \(k\ge0\), (UP5), with \(M>k+1\), permits
differentiating (UP6) under the integral \(k\) times and gives \[
b_r^{(k)}(v)=\frac1{2\pi}\int_{\mathbb R}(iy)^kG_r(iy)e^{iyv}\,dy.
\tag{UP7}
\] These derivatives are continuous by dominated convergence. Moreover,
for \(N\) sufficiently large, Fourier inversion applies to the
integrable Fourier transform of the actual convolution density
\(b_{r,N}\). It gives its continuous representative and the
corresponding derivative formula through any fixed order \(k\). The
uniform domination (UP5) and pointwise convergence imply \[
\sup_{v\in\mathbb R}|b_{r,N}^{(k)}(v)-b_r^{(k)}(v)|
\le\frac1{2\pi}\int_{\mathbb R}|y|^k
|G_{r,N}(iy)-G_r(iy)|\,dy\longrightarrow0.
\tag{UP8}
\] For \(k=0\), every approximating representative is nonnegative and
zero outside \([-r,r]\); its continuous representative has these
properties everywhere since it has them almost everywhere. The uniform
limit therefore has the same properties. Uniform convergence on
\([-r,r]\) gives \(\int b_r=\lim_N\int b_{r,N}=1\). Reflection and
real-valuedness also pass to the limit. We have proved, with its density
and limiting sense specified, \[
b_r\in C_c^\infty(\mathbb R,\mathbb R),\quad b_r\ge0,\quad
b_r(-v)=b_r(v),\quad \int b_r(v)\,dv=1,\quad
\operatorname{supp}b_r\subset[-r,r].
\tag{UP9}
\] This is the fixed infinite convolution meant in (UP1). No random
choice of a bump or dependence on a zero is involved.

Uniform convergence of the densities on their common compact support
permits taking the limit in the integral in (UP2), uniformly for \(z\)
in compact subsets. Consequently \[
G_r(z)=\int_{\mathbb R}b_r(v)e^{-zv}\,dv,
\qquad |G_r(z)|\le e^{r|\Re z|}.
\tag{UP10}
\] The function is real on the real axis, has real Taylor coefficients,
and satisfies \(G_r(-z)=G_r(z)\).

\subsection{2. The exact zero set of the
product}\label{the-exact-zero-set-of-the-product}

The product zero multiplicities in this section also occur in Arias de
Reyna, cited above, equation (9). The coordinate \(\zeta=rz/(2\pi i)\)
carries his nonzero integer zeros to the exact lattice in (UP11),
without changing their orders. The proof below verifies the assertion
directly for the present product.

There are no zeros of (UP3) apart from the zeros of its individual
factors. For proof, at a point where none of the factors vanishes,
choose a small disc on which a finite initial product has no zeros and
the tail deviations from 1 have modulus at most \(1/2\). Their sum
converges uniformly by (UP4). The analytic logarithms of the tail
factors, given by the power series for \(\log(1+w)\), then have a
uniformly absolutely convergent sum on a smaller disc, since
\(|\log(1+w)|\le2|w|\). The tail product equals the exponential of that
sum, and is nonzero. At a point where factors vanish, the same reasoning
applies after removing the finitely many vanishing factors. Only
finitely many can vanish there because \(a_jz\to0\).

The zeros of a single factor are the simple zeros \(z=i\pi n/a_j\), with
\(n\in\mathbb Z\setminus\{0\}\). Their union and the order at each point
are therefore \[
\begin{split}
\{z:G_r(z)=0\}
&=\left\{\frac{2\pi i k}{r}:k\in\mathbb Z\setminus\{0\}\right\}
=\{128\pi i k:k\in\mathbb Z\setminus\{0\}\},\\
\operatorname{ord}_{2\pi i k/r}G_r&=1+\nu_2(|k|).
\end{split}
\tag{UP11}
\] Indeed the \(j\)-th factor vanishes at \(2\pi i k/r\) precisely when
\(2^{j-1}\) divides \(k\), giving the displayed number of factors, each
with order one. In particular, every zero of \(G_r\) is on the imaginary
axis. The construction does not claim that a critical-line zeta zero
cannot lie on that explicitly stated lattice.

\subsection{3. The original primitive boundary
filter}\label{the-original-primitive-boundary-filter}

Retain the BPZ operators and coordinate conventions: \[
T=\frac{d^2}{dv^2}-\frac14,\quad f_r=T b_r,\quad
M_f(s)=\int_{\mathbb R}f(v)e^{-(s-1/2)v}\,dv,\quad
f^\#(v)=\overline{f(-v)}.
\tag{UP12}
\] The test \(f_r\) is real and even and belongs to \(C_c^\infty\), with
support contained in \([-r,r]\). Integration by parts twice has no
boundary terms and proves \[
M_{f_r}(s)
=\bigl((s-\tfrac12)^2-\tfrac14\bigr)G_r(s-\tfrac12)
=s(s-1)G_r(s-\tfrac12).
\tag{UP13}
\] Thus its two endpoint moments vanish, and it is nonzero: \[
M_{f_r}(0)=M_{f_r}(1)=0,
\qquad M_{f_r}(\tfrac12)=\int f_r=-\tfrac14.
\tag{UP14}
\] The factor \(T\) and the original coordinate \(s-\tfrac12\) have not
been rescaled.

Let \(\rho\) be any actual nontrivial zero of \(\xi_R\). BPZ4--BPZ5
gives \(0\le\Re\rho\le1\) and \(\rho\notin\{0,1\}\); its argument does
not need a zero-free assertion on the two vertical boundary lines.
Equations (UP11)--(UP13) give the exact detecting statement \[
\Re\rho\ne\tfrac12\quad\Longrightarrow\quad M_{f_r}(\rho)\ne0.
\tag{UP15}
\] The implication follows because then \(\rho-1/2\) is not purely
imaginary and the factor \(\rho(\rho-1)\) is nonzero. Formula (UP15)
applies simultaneously to all actual off-critical zeros, without
selecting one during the construction of \(b_r\).

\subsection{4. The actual infinite translation
pairing}\label{the-actual-infinite-translation-pairing}

Let \(\rho\) range over the distinct zeros of \(\xi_R\), with their
original multiplicities \(m_\rho\). For compact smooth tests use exactly
the pairing \[
B(f,g)=\sum_\rho m_\rho\overline{M_f(1-\bar\rho)}M_g(\rho),
\quad Q(f)=B(f,f),\quad f_a(v)=f(v-a).
\tag{UP16}
\] The polynomial zero count and compact-test decay proved in BPZ5--BPZ6
ensure absolute convergence. Here is the estimate specialized to the
present test. For every integer \(M\ge0\), integration by parts in the
Fourier factor of the Mellin integral gives \[
\sup_{0\le\sigma\le1}|M_{f_r}(\sigma+i\gamma)|
\le C_M(1+|\gamma|)^{-M}.
\tag{UP17}
\] All derivatives of \(f_r(v)e^{-(\sigma-1/2)v}\) have uniformly
bounded integrals for \(0\le\sigma\le1\), because of (UP9) and its fixed
compact support. These bounds prove (UP17) for \(|\gamma|\ge1\); the
original integral covers the bounded interval. Combining this with BPZ5,
by dyadic height intervals, proves for every integer \(k\ge0\) \[
\sum_\rho m_\rho|M_{f_r}(\rho)|^2(1+|\Im\rho|)^k<\infty.
\tag{UP18}
\]

Since \(f_r\) is real and even, its transform has the exact identities
\[
M_{f_r}(1-s)=M_{f_r}(s),\qquad
M_{f_r}(\bar s)=\overline{M_{f_r}(s)},\qquad
\overline{M_{f_r}(1-\bar\rho)}=M_{f_r}(\rho).
\tag{UP19}
\] Translation in the original variable gives
\(M_{(f_r)_a}(s)=e^{-(s-1/2)a}M_{f_r}(s)\) for every real \(a\).
Therefore, with \(c_\rho=\rho-1/2\), \[
K_r(a):=B((f_r)_a,f_r)
=\sum_\rho m_\rho M_{f_r}(\rho)^2e^{c_\rho a},
\qquad q_r=K_r(0)=Q(f_r).
\tag{UP20}
\] The weight off the critical line is the displayed complex square, not
an absolute square. Every off-critical weight is nonzero by (UP15). The
multiplicity is retained in the coefficient at each distinct exponent.

Since \(|\Re c_\rho|\le1/2\), equations (UP18) and (UP20) justify all
real derivatives locally uniformly: \[
K_r^{(k)}(a)=\sum_\rho m_\rho M_{f_r}(\rho)^2c_\rho^k e^{c_\rho a},
\quad K_r\in C^\infty(\mathbb R),\quad
|K_r(a)|\le C e^{|a|/2}.
\tag{UP21}
\] Changing \(\rho\) to \(1-\rho\), respectively \(\bar\rho\), is
legitimate in these absolutely convergent sums and proves \[
K_r(-a)=K_r(a),\qquad K_r(a)\in\mathbb R
\quad(a\in\mathbb R).
\tag{UP22}
\] No holomorphic extension of this real-time series in the variable
\(a\) is asserted. Its proof of smoothness and the growth bound are the
real-variable statements in (UP21).

\subsection{5. The global Laplace transform and every possible
pole}\label{the-global-laplace-transform-and-every-possible-pole}

For \(\Re w>1/2\), (UP18), the strip bound, and absolute integrability
permit interchanging the integral and full zero sum. They give \[
\mathscr K_r(w):=\int_0^\infty e^{-wa}K_r(a)\,da
=\sum_\rho\frac{m_\rho M_{f_r}(\rho)^2}{w-c_\rho}.
\tag{UP23}
\] For example, the integral of the absolute values of all summands is
bounded by \((\Re w-1/2)^{-1}\sum_\rho m_\rho|M_{f_r}(\rho)|^2\), which
is finite. The series on the right is a meromorphic function on the
entire \(w\)-plane. On a compact set \(|w|\le R\) avoiding its listed
exponent points, only finitely many zeros have bounded height; their
nonzero denominators have a positive minimum. For all sufficiently large
\(|\Im\rho|\), the denominator has modulus at least \(|\Im\rho|/2\), so
(UP18) gives uniform absolute convergence of the remaining tail. The
same argument applies on a disc after removing any one listed term.

It follows that its residue at every possible exponent point is exactly
\[
\operatorname{Res}_{w=c_\rho}\mathscr K_r(w)
=m_\rho M_{f_r}(\rho)^2.
\tag{UP24}
\] The right side is nonzero at every off-critical zero. Distinct actual
zeros have distinct \(c_\rho\), so no sum of other zeros can cancel that
residue. A point with zero weight is a removable singularity, as
specified by (UP24).

Suppose \(K_r\) is bounded on \([0,\infty)\). Its Laplace integral in
(UP23) then defines a holomorphic function on \(\Re w>0\), because on
every compact subset of that half-plane each derivative in \(w\) is
bounded by an integrable multiple of \(a^ke^{-\delta a}\), for some
\(\delta>0\). It equals the meromorphic series for \(\Re w>1/2\). The
half-plane with the discrete exponent points removed is connected: a
segment joining any two points has a compact neighborhood meeting only
finitely many such points, and small arcs detour around them. The
identity theorem consequently extends their equality throughout that
punctured half-plane.

If an actual zero lies off the critical line, reflection provides one
with \(\Re\rho>1/2\). Its exponent \(c_\rho\) lies in that half-plane
and has the nonzero residue (UP24). A holomorphic Laplace integral there
cannot agree on a punctured disc with a meromorphic function having that
residue. This contradiction proves \[
K_r\text{ bounded on }[0,\infty)
\quad\Longrightarrow\quad\mathrm{RH}.
\tag{UP25}
\] This proof uses the whole infinite zero sum and its normally
convergent meromorphic continuation. It does not truncate the actual
spectrum.

Conversely, under RH each \(c_\rho=i\gamma_\rho\) is purely imaginary,
and (UP19) gives \(M_{f_r}(\rho)\in\mathbb R\). Equations (UP18) and
(UP20) then give \[
q_r=\sum_\rho m_\rho M_{f_r}(\rho)^2\ge0,
\qquad |K_r(a)|\le q_r\quad(a\in\mathbb R).
\tag{UP26}
\] This is the exact triangle inequality for that absolutely convergent
sum, with its actual weights. In particular, no independent sign
assumption on \(q_r\) was used to prove (UP25).

\subsection{6. A single fixed test gives the complete
criterion}\label{a-single-fixed-test-gives-the-complete-criterion}

Combining (UP25) and (UP26) proves \[
\boxed{
\begin{split}
\mathrm{RH}
&\Longleftrightarrow K_r\text{ is bounded on }[0,\infty)\\
&\Longleftrightarrow |K_r(a)|\le q_r\text{ for every }a\ge0,
\qquad r=\frac1{64}.
\end{split}}
\tag{UP27}
\] All quantities here arise from the one test fixed in (UP1)--(UP14).
The last displayed inequality implies boundedness directly: at \(a=0\)
it also forces \(q_r\ge0\), and thereafter supplies the finite constant
\(q_r\).

The criterion has an exact two-translate formulation. BPZ13--BPZ14, or
direct cancellation of the exponentials in (UP16), gives \[
Q((f_r)_a)=q_r,
\qquad
Q((f_r)_a+f_r)=2q_r+2K_r(a),
\qquad
Q((f_r)_a-f_r)=2q_r-2K_r(a).
\tag{UP28}
\] Here Hermitian symmetry of \(B\), the reality (UP22), and its
conjugate linearity in the first slot give the cross terms exactly. All
three tests retain both zero endpoint moments because translation
multiplies those moments by exponentials. Consequently (UP27) is also
exactly \[
\mathrm{RH}\quad\Longleftrightarrow\quad
Q((f_r)_a+f_r)\ge0\text{ and }Q((f_r)_a-f_r)\ge0
\text{ for every }a\ge0.
\tag{UP29}
\] If an off-critical zero exists, the proof of (UP25) makes \(K_r\)
unbounded. One can then choose \(a\ge0\) with \(|K_r(a)|>|q_r|\). The
minus sign when \(K_r(a)>0\), and the plus sign when \(K_r(a)<0\), gives
the actual compact primitive test with value \(2q_r-2|K_r(a)|<0\). This
witness requires only a translate and a sign of the same fixed test; it
does not alter the bump in response to the zero.

Continuity also proves an exact countable version: the inequality in
(UP27), or both inequalities in (UP29), can be required for every
rational \(a\ge0\), because every real \(a\ge0\) is a limit of such
rationals. No corresponding criterion on one fixed, equally spaced
translation lattice is asserted here. Sampling on a fixed lattice
identifies distinct imaginary frequencies modulo its period; (UP24),
which separates their continuous Laplace poles, does not justify
discarding that distinction.

\subsection{7. The convolution test and the full supported
receiver}\label{the-convolution-test-and-the-full-supported-receiver}

Define the actual convolution \[
h_r=f_r^\#*f_r=f_r*f_r,
\quad h_r\in C_c^\infty(\mathbb R,\mathbb R),
\quad h_r(-v)=h_r(v),\quad
\operatorname{supp}h_r\subset[-2r,2r].
\tag{UP30}
\] Fubini's theorem on compact supports gives \(M_{h_r}=M_{f_r}^2\). The
convolution in the cross pairing of (UP20) is exactly \[
((f_r)_a)^\#*f_r(v)=h_r(v+a),\qquad
M_{h_r(\,\cdot+a)}(s)=e^{(s-1/2)a}M_{f_r}(s)^2.
\tag{UP31}
\] For proof, \(((f_r)_a)^\#(v)=f_r(v+a)\) because \(f_r\) is real and
even, and a change of variable in the convolution yields the first
identity. Both endpoint values of the transform in (UP31) are zero by
(UP14). Thus the original full explicit formula on this test reads \[
K_r(a)=A_\infty(h_r(\,\cdot+a))
-P_{\rm fin}(h_r(\,\cdot+a)),
\tag{UP32}
\] with the original archimedean and finite-prime terms of BPZ19 and
SZW; no endpoint contribution is omitted. The same statement applies to
each of the two primitive tests in (UP28), using its own full
convolution.

For the finite support semilattice \(\mathscr L\), write
\(\mathbf e_\lambda\) for the retained coordinate basis of the supported
explicit formula. For the test (UP31), its exact vector values are \[
\begin{split}
\boldsymbol B_{\mathscr L}&=0,\\
\boldsymbol Z_{\mathscr L}&=K_r(a)\mathbf e_{1_{\mathscr L}},\\
\boldsymbol D_{\mathscr L}
&=\bigl(P_{\rm fin}-A_\infty\bigr)(h_r(\,\cdot+a))
\mathbf e_{1_{\mathscr L}},\\
\boldsymbol B_{\mathscr L}-\boldsymbol Z_{\mathscr L}
&=\boldsymbol D_{\mathscr L}.
\end{split}
\tag{UP33}
\] Indeed, both top endpoint coefficients and every lower endpoint
coefficient are zero by the evaluated moments in (UP14) and (UP31).
These zeros are values of the existing coordinate functionals. They do
not identify the supported zero element with external absence.

That distinction is retained directly by the synchronized carriers from
BPZ20--BPZ22. For this carrier construction retain the original
nontrivial bounded distributive support lattice \(\mathscr L\), with
bottom \(0_{\mathscr L}\), top \(1_{\mathscr L}\), and
\(0_{\mathscr L}\ne1_{\mathscr L}\). For a complex vector space \(V\),
set \[
G_{\mathscr L}(V)
=\{(0,\lambda):\lambda\in\mathscr L\}
\cup(V\times\{1_{\mathscr L}\}),
\quad A^\uparrow(v,\lambda)=(Av,\lambda)
\tag{UP34}
\] for each complex linear map \(A\). With the original join addition
and meet scalar action, this lift is additive and scalar compatible by
direct substitution. Translation, \(T\), and the endpoint map
\(E(f)=(M_f(0),M_f(1))\) are all such linear maps. Every nonzero
primitive test with label \(1_{\mathscr L}\) maps under \(E^\uparrow\)
to \((0,1_{\mathscr L})\), the supported zero in the endpoint carrier.
The unsupported input \((0,0_{\mathscr L})\) retains its distinct label.
Their entire lower-label families are preserved by (UP34).

Likewise the lifted pairing is exactly \[
B^{\mathscr L}((u,\lambda),(v,\mu))
=(B(u,v),\lambda\wedge\mu).
\tag{UP35}
\] It lies in the synchronized carrier: a nonzero amplitude requires
both input amplitudes nonzero, hence both input labels top. Its
additivity and conjugate scalar compatibility follow from those of \(B\)
and distributivity of meet over join. Thus (UP20) and (UP28) retain top
support even when their amplitudes vanish. The negative amplitude
obtained after (UP29) has this same support label; an endpoint value
equal to supported zero does not change its sign. These formulas keep
all the original label maps while applying the single fixed primitive
test.

\subsection{Reading and scope}\label{reading-and-scope}

The complete programme source read was
\texttt{BOUNDARY\_\allowbreak{}PRIMITIVE\_\allowbreak{}ALL\_\allowbreak{}ZERO\_\allowbreak{}RECEIVER.\allowbreak{}tex}, equations
BPZ1--BPZ23 and all intervening proofs, at the path linked above. Its
original human input is Brad Rodgers and Terence Tao, \emph{The de
Bruijn--Newman constant is non-negative}, arXiv:1801.05914v5, with
source equations \texttt{phidef}, \texttt{htdef}, \texttt{hoz}, and
\texttt{sas}, and the original Weil explicit-formula convention is
retained through the cited programme sources. This note does not
represent a fresh read of those author papers. For the citation
revision, Arias de Reyna's original author TeX \texttt{09-\allowbreak{}Function.\allowbreak{}tex},
arXiv:1702.05442v1, was read at lines 1--720, including Theorem 1 and
equations (1)--(9), their proof, Theorem 4's derivative formula and the
bibliography. The bump and its product zero set are classical results
reproduced here. The programme application is their fixed-test
nonvanishing at every possible off-critical zero and the full BPZ
Laplace argument in the stated coordinates. No arithmetic bound on
\(K_r(a)\) for all translations has been supplied by this proof.

\end{document}
